% SGA 3, Expose VI_A, Loop 1 English reconstruction.
% Scope: Theorem 2.6.5 through Remark 2.6.7;
% Section 3 is excluded.
% Authority: Polo--Gille Exp6A-13oct24.pdf, component-local pages 18--20
% (combined-reader pages 321--323; printed pages 310--312).
% OCR and the Reinhold English Markdown are comparison witnesses only.
% The diagram is an exact 180-dpi crop of the controlling PDF:
% local p. 20, x=650,y=460,w=260,h=170 on a 1488x2105 render.

于是得到下述结果 ([Per75], II, 定理~\SGARefLink{prop:sga3-VIA-2.4}{2.4};
亦见 [Per76], 命题~4.1.1).

\medskip
\phantomsection\label{thm:sga3-VIA-2.6.5}%
\noindent\textbf{定理 2.6.5.}
\textit{设 \(k\) 为域, \(G\) 为 \(k\)-群概形.}
\begin{enumerate}
\item[\SGAInlineRefTarget{thm:sga3-VIA-2.6.5:item:i}(i)] \textit{\(G\) 有唯一的子群概形 \(G^0\), 称为 \(G\) 的\emph{单位分支}, 满足:}
  \begin{enumerate}
  \item[\SGAInlineRefTarget{thm:sga3-VIA-2.6.5:item:a}(a)] \textit{\(G^0\) 的底空间是单位元所在的不可约分支;}
  \item[\SGAInlineRefTarget{thm:sga3-VIA-2.6.5:item:b}(b)] \textit{\(G^0\to G\) 是平坦闭浸入; 即
  \[
    \mathcal O_{G,g}=\mathcal O_{G^0,g}
    \qquad\text{对每个 }g\in G^0.
  \]}
  \end{enumerate}
\item[\SGAInlineRefTarget{thm:sga3-VIA-2.6.5:item:ii}(ii)] \textit{此外, \(G^0\) 拟紧、几何不可约, 且为 \(G\) 的特征子群.}
\item[\SGAInlineRefTarget{thm:sga3-VIA-2.6.5:item:iii}(iii)] \textit{若 \(G\) 连通, 则 \(G=G^0\).}
\end{enumerate}

\noindent\textit{证明.}
\SGARefLink{thm:sga3-VIA-2.6.5:item:i}{(i)} 首先回顾, \(G\) 分离 (\hyperref[sga3:VIA:0.3]{0.3}), 从而当然拟分离.
设 \(U\) 为 \(G\) 的包含 \(C^0\) 的泛点 \(\omega\) 的仿射开子集.
由\hyperref[rem:sga3-VIA-2.6.3]{2.6.3} 和\hyperref[prop:sga3-VIA-2.6.4]{2.6.4},
\[
  \overline U=C^0U
\]
既开又闭, 且在集合意义下,
\[
  C^0=\overline{\{\omega\}}
\]
为这些既开又闭的子集之交.

当 \(U\) 遍历包含 \(\omega\) 的仿射开子集时,
\(G\)-概形 \(\overline U\) 构成一个逆系,
其过渡态射为闭浸入, 因而仿射.
所以可以构造其逆极限 \(G^0\) (参见 EGA~IV\textsubscript{3}, 8.2.2).
因此, 对 \(G\) 的每个仿射开集 \(V\), \(G^0\cap V\) 是下列代数的谱:
\[
  \varinjlim_U\mathcal O_G(V\cap\overline U)
  =
  \mathcal O_G(V)\Big/\sum_U I_{\overline U}(V),
\]
其中 \(I_{\overline U}(V)\) 表示下述态射的核:
\[
  \mathcal O_G(V)\longrightarrow
  \mathcal O_G(V\cap\overline U).
\]
从而 \(G^0\) 的底空间是 \(C^0\), 且 \(G^0\to G\) 是闭浸入.

此外, 对每个 \(g\in G^0\), \(\mathcal O_{G^0,g}\) 是下列
\(k\)-代数的归纳极限, 其中 \(V\) 遍历 \(G\) 的包含 \(g\) 的仿射开子集:
\[
  \mathcal O_{G^0}(V\cap G^0)
  =
  \varinjlim_U\mathcal O_G(V\cap\overline U).
\]
此二重归纳极限可等同于
\[
  \varinjlim_U\
  \varinjlim_{\substack{V\\g\in V\subset\overline U}}
  \mathcal O_G(V)
  =
  \mathcal O_{G,g}.
\]
因此
\[
  \mathcal O_{G^0,g}=\mathcal O_{G,g}
\]
(亦见 EGA~IV\textsubscript{2}, 5.13.3(ii)), 因而 \(i:G^0\to G\) 是平坦闭浸入.
反之, 此条件蕴涵 \(i^*(\mathcal O_G)=\mathcal O_{G^0}\),
故 \(G^0\) 由条件\SGARefLink{thm:sga3-VIA-2.6.5:item:a}{(a)}
和\SGARefLink{thm:sga3-VIA-2.6.5:item:b}{(b)} 唯一确定.
这就证明了\SGARefLink{thm:sga3-VIA-2.6.5:item:i}{(i)}.

\SGARefLink{thm:sga3-VIA-2.6.5:item:ii}{(ii)} 的前两个断言由
\hyperref[cor:sga3-VIA-2.6.2]{2.6.2} 得到.
最后, 设 \(S\) 为 \(k\)-概形, \(\phi\) 为 \(S\)-群 \(G_S\) 的自同构.
对每个 \(s\in S\), \(\phi_s\) 将 \(G^0_{\kappa(s)}\) 映入自身,
故 \(\phi(G^0_S)\subset G^0_S\). 此外, 由 \(i\) 经基变换得到的闭浸入
\[
  i_S:G^0_S\hookrightarrow G_S
\]
平坦, 故
\[
  \mathcal O_{G_S,\phi(z)}
  =
  \mathcal O_{G^0_S,\phi(z)}
  \qquad\text{对每个 }z\in G^0_S.
\]
因此 \(\phi\circ i_S\) 经由 \(G^0_S\) 分解.
这就证明了 \(G^0\) 是 \(G\) 的特征子群, 从而证明了\SGARefLink{thm:sga3-VIA-2.6.5:item:ii}{(ii)}.
最后, \SGARefLink{thm:sga3-VIA-2.6.5:item:iii}{(iii)}
是下述命题的\SGARefLink{thm:sga3-VIA-2.6.5:item:i}{(i)} 的特殊情形.%
\footnote{\label{note:sga3-VIA-32}编者注: 此命题以及前面的结果均由 O.~Gabber 告知我们.
它将用于修正第 VI\textsubscript{B} 篇\hyperref[thm:sga3-VIB-5.3]{定理~5.3} 的证明.}

\medskip
\phantomsection\label{prop:sga3-VIA-2.6.6}%
\noindent\textbf{命题 2.6.6.}
\textit{设 \(k\) 为域, \(G\) 为作用于 \(k\)-概形 \(X\) 的 \(k\)-群, 使得
\[
  G\times X\longrightarrow X\times X,\qquad
  (g,x)\longmapsto(gx,x)
\]
为满射. 假设 \(X\) 拟分离. 则:}
\begin{enumerate}
\item[\SGAInlineRefTarget{prop:sga3-VIA-2.6.6:item:i}(i)] \textit{\(X\) 的每个连通分支 \(C\) 都不可约.}
\item[\SGAInlineRefTarget{prop:sga3-VIA-2.6.6:item:ii}(ii)] \textit{设 \(\eta\) 为 \(C\) 的泛点, \(L\) 为 \(k\) 在 \(\kappa(\eta)\) 中的代数闭包.
则 \(C\) 为 \(L\)-概形, 在 \(L\) 上几何不可约, 且态射
\[
  G^0_L\times_LC\longrightarrow C\times_LC
\]
为满射.}
\item[\SGAInlineRefTarget{prop:sga3-VIA-2.6.6:item:iii}(iii)] \textit{特别地, 若 \(C\) 包含有理点 \(x\),
则 \(C\) 在 \(k\) 上几何不可约, 且
\[
  \phi:G^0\longrightarrow C,\qquad g\longmapsto gx,
\]
为满射.}
\end{enumerate}

\noindent\textit{证明.}
\SGARefLink{prop:sga3-VIA-2.6.6:item:i}{(i)} 设 \(C\) 为 \(X\) 的连通分支, \(Y\) 为 \(X\) 中包含于 \(C\) 的不可约分支,
\(\eta\) 为 \(Y\) 的泛点, \(U\) 为 \(X\) 的包含 \(\eta\) 的仿射开子集.
由于 \(X\) 拟分离, \(U\) 在 \(X\) 中相对拟紧.
故由\hyperref[prop:sga3-VIA-2.6.4]{2.6.4},
\[
  \overline U=G^0U
\]
为 \(X\) 的既开又闭的子集, 与 \(C\) 相交, 从而包含 \(C\).
由\hyperref[rem:sga3-VIA-2.6.3]{2.6.3}, 当 \(U\) 遍历 \(\eta\)
的一个仿射开邻域基时, 各 \(\overline U\) 的交为 \(\overline{\{\eta\}}\).
从而
\[
  C=\overline{\{\eta\}},
\]
这就证明了\SGARefLink{prop:sga3-VIA-2.6.6:item:i}{(i)}.

\SGARefLink{prop:sga3-VIA-2.6.6:item:ii}{(ii)} 和\SGARefLink{prop:sga3-VIA-2.6.6:item:iii}{(iii)}
的第一个断言均由\hyperref[prop:sga3-VIA-2.6.1]{2.6.1} 得到.
先证明\SGARefLink{prop:sga3-VIA-2.6.1:item:iii}{(iii)} 的第二个断言.
用 \(\omega\), 相应地 \(\eta\), 表示 \(G^0\), 相应地 \(C\), 的泛点.
设 \(z\in C\), 并令 \(K=\kappa(z)\).
由于 \(G^0\), 相应地 \(C\), 在 \(k\) 上几何不可约,
概形 \(G^0_K\), 相应地 \(C_K\), 不可约;
用 \(\xi\), 相应地 \(\beta\), 表示其泛点.
在\hyperref[prop:sga3-VIA-2.6.4]{2.6.4} 的证明中已见,
\[
  \mu_z:G_K\longrightarrow C_K,\qquad h\longmapsto h\cdot z,
\]
将 \(\xi\) 映到 \(\beta\). 同样, \(\phi(\omega)=\eta\).

令 \(L=\kappa(\xi)\), 并设 \(\xi_L\) 和 \(z_L\)
为由 \(\xi\) 及 \(z\) 得到的 \(L\)-点. 则
\[
  \xi_L\cdot z_L=\beta'
\]
为 \(C_L\) 中位于 \(\beta\in C_K\) 上方的一点,
因而也位于 \(\eta\in C\) 上方. 考虑笛卡儿方图
\begin{center}
  \phantomsection\label{diag:sga3-VIA-2.6.6-cartesian-square}%
  \begin{tikzcd}[row sep=large,column sep=large]
    G_L^0 \arrow[r,"\varphi_L"] \arrow[d,"\pi_{G^0}"'] &
    C_L \arrow[d,"\pi_C"] \\
    G^0 \arrow[r,"\varphi"'] &
    C
  \end{tikzcd}
\end{center}
由于
\[
  \phi_L\bigl(\pi_{G^0}^{-1}(\omega)\bigr)
  =
  \pi_C^{-1}(\eta)
\]
(参见 EGA~I, 3.4.8), 存在 \(g\in G^0_L\),
使 \(\phi_L(g)=\beta'\). 因此
\[
  z_L
  =
  \xi_L^{-1}\cdot\phi_L(g)
  =
  \phi_L(\xi_L^{-1}g),
\]
从而
\[
  \phi\bigl(\pi_{G^0}(\xi_L^{-1}g)\bigr)
  =
  \pi_C(z_L)
  =
  z.
\]
这就证明了 \(\phi\) 满射.

现在证明\SGARefLink{prop:sga3-VIA-2.6.6:item:ii}{(ii)} 的第二个断言.
只需证明对每个 \(z\in C\), 态射
\[
  \mu_z:
  G^0_L\otimes_L\kappa(z)
  \longrightarrow
  C\otimes_L\kappa(z)
\]
为满射. 但这由\SGARefLink{prop:sga3-VIA-2.6.6:item:iii}{(iii)} 得到,
因为 \(z\) 为 \(C\otimes_L\kappa(z)\) 的有理点.

\medskip
\phantomsection\label{rem:sga3-VIA-2.6.7}%
\noindent\textbf{注 2.6.7.---}
\textit{在\hyperref[prop:sga3-VIA-2.6.6]{2.6.6} 的假设下,
若 \(C(k)=\varnothing\), 则态射
\[
  G^0\times_kC\longrightarrow C\times_kC
\]
未必满射. 例如, 设 \(k=\mathbb R\), \(G=\{\pm1\}_{\mathbb R}\).
\(G\)-挠子
\[
  X=\operatorname{Spec}\mathbb R[X]/(X^2+1)
\]
连通, 但
\[
  G^0\times_{\mathbb R}X\longrightarrow X\times_{\mathbb R}X
\]
并不满射. (这里 \(L=\mathbb C\), 且
\[
  G^0\times_{\mathbb R}X\longrightarrow X\times_{\mathbb C}X
\]
为同构.)}
